%% kapa: exam=1 round=22 year=2011 period=2 form=A subject=accounting subject_ko=회계학 seq=8 number=8
%% source: Q-Net 감정평가사 기출문제 (cst00309) artlSeq=5036790 file=extracted/1008/5036790/2011년 제22회 감정평가사 1차시험 기출문제/2교시.hwp
%% conversion: kordoc 4.24.51 -> tools/gen1.py (auto); review=proofread-gpt-astra; flags=none
\begin{question}{8}
용역의 제공에 따른 수익인식기준에 관한 설명으로 옳지 않은 것은?
\begin{choices}
  \item 광고제작수수료는 진행기준에 따라 수익을 인식한다.
  \item 하나의 공연입장권으로 여러 행사에 참여할 수 있는 경우에는 각각 행사를 위한 용역 수행정도에 따라 각 행사에 배분하여 수익을 인식한다.
  \item 학원수강료는 강의기간에 걸쳐 수익으로 인식한다.
  \item 창업운영지원용역의 제공에 대한 수수료는 수취시점에서 수익으로 인식한다.
  \item 회원이 회비를 납부하고 회원가입기간 동안 재화나 용역을 비회원보다 저렴한 가격으로 구매할 수 있는 경우에는 이 효익이 제공되는 시기, 성격 및 가치를 반영하는 기준으로 수익을 인식한다.
\end{choices}
\end{question}
\answer{4}
